\documentclass[preview]{standalone}

\usepackage{amsmath}
\usepackage{amssymb}
\usepackage{stellar}
\usepackage{definitions}

\begin{document}

\id{algebras-over-fields-exercises}
\genpage

\section{Exercises}

\begin{snippetexercise}{rings-ex3}{A matrix representation of the quaternions}
    Consider
    \[
        A=
        \left\{
            \begin{pmatrix}
                z&w\\
                -\overline w&\overline z
            \end{pmatrix}
            \suchthat z,w\in\complexnumbers
        \right\}
        \subseteq\matrices_2(\complexnumbers).
    \]
    Show that \(A\) is isomorphic, as a real algebra, to the quaternion
    algebra
    \[
        \mathbb H
        =
        \{a+b\mathbf i+c\mathbf j+d\mathbf k
          \suchthat a,b,c,d\in\realnumbers\}.
    \]
\end{snippetexercise}

\begin{snippetsolution}{rings-ex3-sol}{A matrix representation of the quaternions}
    Writing \(z=a+bi\) and \(w=c+di\), every matrix in \(A\) has the
    unique expression
    \[
        aE_0+bE_1+cE_2+dE_3,
    \]
    where
    \[
        E_0=\begin{pmatrix}1&0\\0&1\end{pmatrix},
        \qquad
        E_1=\begin{pmatrix}i&0\\0&-i\end{pmatrix},
    \]
    \[
        E_2=\begin{pmatrix}0&1\\-1&0\end{pmatrix},
        \qquad
        E_3=\begin{pmatrix}0&i\\i&0\end{pmatrix}.
    \]
    Thus \(E_0,E_1,E_2,E_3\) form a real basis of \(A\).

    Direct multiplication gives
    \[
        E_1^2=E_2^2=E_3^2=-E_0,
    \]
    \[
        E_1E_2=E_3,
        \qquad
        E_2E_3=E_1,
        \qquad
        E_3E_1=E_2,
    \]
    while reversing the order changes the sign. These are exactly the
    defining multiplication relations of
    \(1,\mathbf i,\mathbf j,\mathbf k\).

    Therefore the \(\realnumbers\)-linear \function
    \[
        \Phi\colon\mathbb H\fromto A,
        \qquad
        \Phi(a+b\mathbf i+c\mathbf j+d\mathbf k)
        =
        \begin{pmatrix}
            a+bi&c+di\\
            -c+di&a-bi
        \end{pmatrix}
    \]
    maps a basis to a basis and preserves multiplication. It is an
    isomorphism of real algebras.
\end{snippetsolution}

\begin{snippetnote}{quaternion-matrix-representation-real-algebra-note}{The scalar field}
    The isomorphism is one of \(\realnumbers\)-algebras, not of
    \(\complexnumbers\)-algebras. The usual quaternion algebra has no
    canonical compatible complex scalar multiplication.
\end{snippetnote}

\end{document}
